% from ega0-1-fr.tex
\oldpage[0\textsubscript{I}]{64}

  \item[$b)$] L'homomorphisme $v$ est bicontinu~;
  \item[$c)$] $A$ est un anneau $\mathfrak J$-adique.
\end{enumerate}
\end{corollary}

\begin{corollary}[7.2.6]
\label{0.7.2.6-fr}
Soient $A$ un anneau admissible, $\mathfrak J$ un idéal de définition de $A$.
Pour que $A$ soit noethérien, il faut et il suffit que $A/\mathfrak J$ soit
noethérien et que $\mathfrak J/\mathfrak J^2$ soit un
$(A/\mathfrak J)$-module de type fini.
\end{corollary}

Ces conditions sont évidemment nécessaires. Inversement, supposons-les
vérifiées~; comme en vertu de (7.2.4) $A$ est complet pour la topologie
$\mathfrak J$-préadique, pour qu'il soit noethérien, il faut et il suffit que
l'anneau gradué associé $\operatorname{grad}(A)$ (pour la filtration des
$\mathfrak J^n$) le soit ([I], p.~18-07, th.~4). Or, soient
$a_1,\ldots,a_n$ des éléments de $\mathfrak J$ dont les classes mod.
$\mathfrak J^2$ sont des générateurs de $\mathfrak J/\mathfrak J^2$ en tant
que $A/\mathfrak J$-module. Il est immédiat par récurrence que les classes mod.
$\mathfrak J^{m+1}$ des monômes de degré total $m$ en les $a_i$
$(1\leq i\leq n)$ forment un système de générateurs du
$(A/\mathfrak J)$-module $\mathfrak J^m/\mathfrak J^{m+1}$. On en conclut que
$\operatorname{grad}(A)$ est un anneau isomorphe à un quotient de
$(A/\mathfrak J)[T_1,\ldots,T_n]$ ($T_i$ indéterminées), ce qui achève la
démonstration.

\begin{proposition}[7.2.7]
\label{0.7.2.7-fr}
Soit $(A_i,u_{ij})$ un système projectif $(i\in\mathbf N)$ d'anneaux discrets,
et pour tout entier $i$, soit $\mathfrak J_i$ le noyau dans $A_i$ de
l'homomorphisme $u_{0i}:A_i\to A_0$. On suppose que~:
\begin{enumerate}
  \item[$a)$] Pour $i\leq j$, $u_{ij}$ est surjectif et son noyau est
    $\mathfrak J_j^{i+1}$ (donc $A_i$ est isomorphe à
    $A_j/\mathfrak J_j^{i+1}$).
  \item[$b)$] $\mathfrak J_1/\mathfrak J_1^2$ ($=\mathfrak J_1$) est un
    module de type fini sur $A_0=A_1/\mathfrak J_1$.
\end{enumerate}
Soit $A=\varprojlim_i A_i$, et pour tout entier $n\geq0$, soient $u_n$
l'homomorphisme canonique $A\to A_n$, $\mathfrak J^{(n+1)}\subset A$ son
noyau. Dans ces conditions~:
\begin{enumerate}
  \item[(i)] $A$ est un anneau adique, ayant pour idéal de définition
    $\mathfrak J=\mathfrak J^{(1)}$.
  \item[(ii)] On a $\mathfrak J^{(n)}=\mathfrak J^n$ pour tout $n\geq1$.
  \item[(iii)] $\mathfrak J/\mathfrak J^2$ est isomorphe à
    $\mathfrak J_1=\mathfrak J_1/\mathfrak J_1^2$, et est par suite un module
    de type fini sur $A_0=A/\mathfrak J$.
\end{enumerate}
\end{proposition}

L'hypothèse de surjectivité des $u_{ij}$ entraîne que $u_n$ est surjectif~;
en outre, l'hypothèse $a)$ implique que $\mathfrak J_j^{j+1}=0$, donc $A$ est
un anneau admissible (7.2.2)~; par définition, les $\mathfrak J^{(n)}$ forment
un système fondamental de voisinages de $0$ dans $A$, donc (ii) entraîne (i).
En outre, on a $\mathfrak J=\varprojlim_i\mathfrak J_i$ et les applications
$\mathfrak J\to\mathfrak J_i$ sont surjectives, donc (ii) entraîne (iii), et
on est ramené à prouver (ii). Par définition, $\mathfrak J^{(n)}$ est formé
des éléments $(x_k)_{k\geq0}$ de $A$ tels que $x_k=0$ pour $k<n$, donc
$\mathfrak J^{(n)}\mathfrak J^{(m)}\subset\mathfrak J^{(n+m)}$, autrement dit
les $\mathfrak J^{(n)}$ constituent une \emph{filtration} de $A$. D'autre
part, $\mathfrak J^{(n)}/\mathfrak J^{(n+1)}$ est isomorphe à la projection de
$\mathfrak J^{(n)}$ sur $A_n$~; comme
$\mathfrak J^{(n)}=\varprojlim_{i>n}\mathfrak J_i^n$, cette projection n'est
autre que $\mathfrak J_n^n$, qui est un module sur
$A_0=A_n/\mathfrak J_n$. Soient alors $a_j=(a_{jk})_{k\geq0}$ $r$ éléments de
$\mathfrak J=\mathfrak J^{(1)}$ tels que $a_{11},\ldots,a_{r1}$ forment un
système de générateurs de $\mathfrak J_1$ sur $A_0$~; nous allons voir que
l'ensemble $S_n$ des monômes de degré total $n$ en les $a_j$ engendre l'idéal
$\mathfrak J^{(n)}$ de $A$. Comme $\mathfrak J_i^{i+1}=0$, il est clair tout
d'abord que $S_n\subset\mathfrak J^{(n)}$~; puisque $A$ est complet pour la
filtration $(\mathfrak J^{(m)})$, il suffit de prouver que l'ensemble
$\overline S_n$ des classes mod. $\mathfrak J^{(n+1)}$ des éléments de $S_n$
engendre le module gradué $\operatorname{grad}(\mathfrak J^{(n)})$ sur
l'anneau gradué $\operatorname{grad}(A)$ pour la filtration précédente
([I], p.~18-06, lemme)~; en vertu de la définition de la multiplication dans
$\operatorname{grad}(A)$,

